\section{Conclusion}

\label{sec:conclusions}
We introduced \emph{Neural LoFi}, a tractable spectral theory of multilayer feature learning beyond the lazy regime. The framework turns deep representation learning into an explicit layerwise procedure: given a current representation, each layer selects low-degree features that are both task-correlated and simple in the geometry induced by previous layers. This yields a concrete relevance--complexity principle for feature selection, a quantitative criterion for when new features emerge, and a principle of \emph{low-degree compositionality}: depth helps when each layer makes the next useful signal statistically detectable.

This picture is visible beyond idealized models. Neural LoFi improves over random-feature baselines, predicts the sample scale at which layerwise features emerge, aligns with early gradient-descent representations, and recovers salient convolutional filters without backpropagation. Together, these results suggest that part of deep feature learning is already captured by supervised low-degree spectral structure.

The deliberate simplicity of Neural LoFi also delineates a clear research program. Theoretically, it remains to understand the long-time dynamics of the task-adaptive kernels generated by LoFi, their behavior beyond second-order filtering, and their extension to structured architectures such as attention. 

Algorithmically, the main challenge is to scale LoFi from a mechanistic surrogate into a useful training primitive---for initialization, feature pretraining, pruning, or diagnostics. Incorporating data reuse, multi-pass dynamics, and backward feature correction would bring the framework closer to trained networks.
